% ===========================================================================
% VERSIÓN OBSOLETA — NO ENVIAR A REVISTA
%
% Esta es la versión INGLESA ANTIGUA del paper. Contiene errores corregidos
% en la versión española (dark_forest_paper_es.tex):
%   - Ec.(3) cuadrática (P_det < 0 para d > r_eff) → lorentziana en ES
%   - fp=0.5, ne=2 → fp=1.0, ne=0.22 en ES
%   - α=0.03 yr⁻¹ → α=0.001 yr⁻¹ en ES
%   - Afiliación USACH → Investigador independiente en ES
%   - Sin simulaciones de control → §2.6 añadida en ES
%
% Si se requiere versión inglesa: regenerar desde dark_forest_paper_es.tex
% ===========================================================================

\documentclass[useAMS,usenatbib]{mn2e}

% ─── PACKAGES ────────────────────────────────────────────────────────────────
\usepackage{amsmath,amssymb}
\usepackage{graphicx}
\usepackage{booktabs}
\usepackage{array}
\usepackage{multirow}
\usepackage{xcolor}
\usepackage{hyperref}
\usepackage{url}
\usepackage{siunitx}
\usepackage{textcomp}

% \doi not defined by mn2e — provide a hyperlinked definition
\newcommand{\doi}[1]{\href{https://doi.org/#1}{\nolinkurl{doi:#1}}}

% ─── JOURNAL METADATA ────────────────────────────────────────────────────────
\title[Dark Forest Monte Carlo]{
  A Monte Carlo Model of the Dark Forest Hypothesis:\\
  Stochastic Detection--Destruction Dynamics and\\
  Spatial Clustering Inversion Across Galactic Population Regimes
}

\author[J.~Galaz]{
  Juan Galaz$^{1}$\thanks{E-mail: juan.galaz.a@usach.cl;
  ORCID: 0009-0007-7474-7560}
  \\
  $^{1}$Departamento de Ingenier\'ia en Minas,
  Universidad de Santiago de Chile (USACH),
  Av. Libertador Bernardo O'Higgins 3363,
  Estaci\'on Central, Santiago, Chile
}

\date{Accepted 2026. Received 2026; in original form 2026}

\pubyear{2026}

\begin{document}

\maketitle

% ─── ABSTRACT ────────────────────────────────────────────────────────────────
\begin{abstract}
We present a stochastic Monte Carlo simulation of the Dark Forest
hypothesis, the conjecture that civilizations in the galaxy preemptively
destroy one another upon mutual detection, and characterise the
resulting destruction rates and spatial distributions across two
population regimes.
Under \emph{pessimistic} Drake parameters ($N_c = 440$, 10 independent
seeds) the model yields a destruction rate of $16.1 \pm 1.8$ per cent,
with surviving civilizations distributed at a mean nearest-neighbour
distance of $1842 \pm 60$\,ly; a Kolmogorov--Smirnov test confirms
that this distribution is significantly super-Poissonian
($p = 1.91\times10^{-42}$, $\text{ratio}_\text{obs/Poisson} = 1.42 \pm 0.05$),
indicating that Dark Forest dynamics create spatial voids around survivors.
Under \emph{optimistic} Drake parameters ($N_c = 10{,}000$ cap, 10 seeds)
the destruction rate rises to $85.8 \pm 0.2$ per cent, with survivors
concentrated at $478 \pm 17$\,ly;
critically, the obs/Poisson ratio \emph{inverts} to $0.578 \pm 0.020$
(sub-Poissonian), indicating clustering rather than dispersal:
a regime transition not previously reported in the literature.
A global Sobol sensitivity analysis ($N = 512$, 5{,}120 evaluations)
identifies no single dominant parameter, with first-order indices
$S_1 \in [0.09, 0.12]$ for all four parameters examined, confirming
that the qualitative results are robust to parameter uncertainty.
These findings suggest a density-dependent inversion in the spatial
structure of surviving civilizations: sparse galaxies produce
super-Poissonian survivors (Dark Forest voids), while dense galaxies
produce sub-Poissonian survivors (forced clustering near GHZ refugia).
\end{abstract}

\begin{keywords}
astrobiology --- extraterrestrial intelligence --- Fermi paradox ---
Monte Carlo methods --- spatial statistics --- galactic habitable zone
\end{keywords}

% ─────────────────────────────────────────────────────────────────────────────
\section{Introduction}
\label{sec:intro}
% ─────────────────────────────────────────────────────────────────────────────

The Fermi paradox, the apparent contradiction between the high
prior probability of extraterrestrial technological civilizations and
their observational absence, has motivated a wide class of
explanatory hypotheses
\citep{fermi1950,hart1975,brin1983,webb2002}.
Among the most radical is the \emph{Dark Forest hypothesis} (DFH),
in which the silence of the cosmos is a strategic equilibrium: any
civilization that reveals its location risks annihilation by a more
advanced neighbour, and therefore the optimal strategy is either
silence or preemptive attack
\citep{cixin2008,devladar2013}.

Formal models of the DFH remain scarce in the peer-reviewed
astrobiological literature.
\citet{forgan2009} developed the most widely cited numerical testbed
for Fermi-paradox hypotheses, using Monte Carlo realisations of the
Drake equation, but modelled destruction as a fixed probability
($P_\text{destroy} = 0.5$) independent of detection, distance or
technological disparity.
\citet{devladar2013} formalised the game-theoretic foundation of the
DFH, showing that preemptive attack is a Nash equilibrium under
certain payoff matrices, but did not simulate the galactic-scale
spatial consequences.
\citet{foucher2023} and \citet{forgan2011} independently noted that
the spatial distribution of survivors should depend on population
density, but neither quantified the regime transition.

The present work addresses this gap by implementing a full
detection--decision--destruction Monte Carlo pipeline with the
following novel elements:
\begin{enumerate}
  \item A continuous three-dimensional detection probability function
        based on technological age and a physically motivated maximum
        detection radius \citep[$r_\text{max} = 1000$\,ly,][]{troitskij1989,
        loeb2007};
  \item Heterogeneous sociological types (aggressive, defensive,
        pacifist, curious) with type-dependent attack thresholds
        following \citet{devladar2013};
  \item A continuous Galactic Habitable Zone (GHZ) density profile
        based on \citet{lineweaver2004}, replacing the uniform-ring
        approximation;
  \item A historical memory of destroyed civilizations that reduces
        communication probability in hazardous directions;
  \item Dual-scenario analysis (pessimistic/optimistic Drake
        parameters) enabling regime comparison;
  \item Global Sobol sensitivity analysis with $N = 512$
        Saltelli samples \citep{saltelli2010}.
\end{enumerate}

The central result, a density-dependent inversion of the
obs/Poisson spatial ratio, has not been reported previously and
constitutes the primary scientific contribution of this paper.

% ─────────────────────────────────────────────────────────────────────────────
\section{Model}
\label{sec:model}
% ─────────────────────────────────────────────────────────────────────────────

\subsection{Civilisation generation}
\label{sec:civgen}

Civilisations are generated by sampling the Drake equation:
\begin{equation}
  N_c = N_\star \cdot f_p \cdot n_e \cdot f_l \cdot f_i \cdot f_c
  \label{eq:drake}
\end{equation}
where $N_\star = 2\times10^{11}$ is the number of Milky Way stars
suitable for life-bearing planets, $f_p = 0.5$ the fraction with
planets, $n_e = 2$ the mean number of habitable-zone planets per
system, and $f_l$, $f_i$, $f_c$ are the life, intelligence and
communication fractions whose product defines the scenario (Table
\ref{tab:scenarios}).

Galactic positions are drawn from the continuous GHZ density profile
\citep{lineweaver2004}:
\begin{equation}
  f(r) \propto \exp\!\left[-\frac{(r - r_\text{peak})^2}{2\sigma_r^2}\right],
  \quad r_\text{peak} = 25{,}000\,\text{ly},\quad \sigma_r = 8{,}000\,\text{ly}
  \label{eq:ghz}
\end{equation}
with a disc half-height of 300\,ly sampled uniformly.
Birth times are drawn uniformly over the galaxy's estimated age of
$1.36\times10^{10}$\,yr, and technological lifetimes follow a
log-normal distribution with median $10^4$\,yr
\citep{sandberg2018}.

\subsection{Sociological heterogeneity}
\label{sec:socio}

Each civilisation is assigned one of four sociological types:
\emph{aggressive} (probability 0.30), \emph{defensive} (0.40),
\emph{pacifist} (0.20), \emph{curious} (0.10).
The type determines the attack-threshold parameter
$\theta_\text{threat}$ and the response to neighbourhood context.
An aggressive civilisation attacks whenever the relative
technological disparity of a detected neighbour exceeds $\theta_\text{threat}$;
a pacifist civilisation prioritises concealment; a curious civilisation
increases its emission radius by a factor of three but reduces it by
80 per cent within 2{,}000\,ly of a previously destroyed civilisation
(historical memory of hazard zones).

\subsection{Detection probability}
\label{sec:detection}

The probability of civilisation $i$ detecting civilisation $j$ at
Euclidean distance $d$ at time $t$ is:
\begin{equation}
  P_\text{det}(d, \tau_i, \nu_i) =
    \left[1 - \left(\frac{d}{r_\text{eff}(\tau_i, \nu_i)}\right)^2\right]
    \cdot (1 - p_\text{fn}) + p_\text{fp}
  \label{eq:pdet}
\end{equation}
where $\tau_i$ is the technological age of the hunter, $\nu_i$ its
technological level, $p_\text{fn} = 0.05$ the false-negative rate, and
$p_\text{fp} = 0.01$ the false-positive rate.
The effective detection radius grows logistically:
\begin{equation}
  r_\text{eff}(\tau, \nu) =
    \frac{r_\text{max}}{1 + (r_\text{max}-1)\,e^{-\alpha\tau}}
    \cdot (1 + \max(0,\nu))
  \label{eq:reff}
\end{equation}
with $r_\text{max} = 1000$\,ly \citep{troitskij1989,loeb2007} and
growth rate $\alpha = 0.03$\,yr$^{-1}$.
A 3-D cKDTree spatial index \citep{maneewongvatana1999} limits
neighbour queries to civilisations within the light-travel horizon,
reducing the per-event complexity from $O(N^2)$ to $O(N \log N)$.

\subsection{Extended lifetime}
\label{sec:lifetime}

To allow interactions beyond the mean technological lifetime $L$,
civilisations remain detectable for an extended window
$L_\text{ext} = 1.43\times10^8$\,yr after birth \citep{sandberg2018}.
This parameter ensures that the simulation samples a sufficient
number of coexisting civilisation pairs across all seeds.

\subsection{Sensitivity analysis}
\label{sec:sobol}

Global sensitivity indices \citep{saltelli2010} were computed for
four parameters: $\theta_\text{threat} \in [0.3,0.9]$,
$\alpha \in [0.001,0.1]$\,yr$^{-1}$, $f_\text{aggressive} \in [0.1,0.5]$,
and $r_\text{max} \in [50,500]$\,ly.
Saltelli sampling with $N = 512$ base samples yields 5{,}120 model
evaluations; results are reported in Section \ref{sec:sobol_results}.

% ─────────────────────────────────────────────────────────────────────────────
\section{Results}
\label{sec:results}
% ─────────────────────────────────────────────────────────────────────────────

\subsection{Destruction rates}
\label{sec:destruction}

Table \ref{tab:results} summarises the destruction statistics for
both scenarios across 10 independent seeds.

\begin{table*}
  \centering
  \caption{Summary statistics for pessimistic and optimistic scenarios
           (10 seeds each). Uncertainties are $\pm 1\sigma$ across seeds.}
  \label{tab:results}
  \begin{tabular}{lcccccc}
    \toprule
    Scenario & $N_c$ & Destroyed & Rate (\%) &
    $\bar{d}_\text{nn}$ (ly) & $r_\text{obs/P}$ & KS $p$-value \\
    \midrule
    Pessimistic & 440 & $70.6 \pm 8.0$ & $16.1 \pm 1.8$ &
    $1842 \pm 60$ & $1.42 \pm 0.05$ & $1.91\times10^{-42}$ \\
    Optimistic  & $10{,}000^\dagger$ & $8583 \pm 22$ & $85.8 \pm 0.2$ &
    $478 \pm 17$  & $0.578 \pm 0.020$ & $<2.2\times10^{-308\ddagger}$ \\
    \bottomrule
  \end{tabular}
  \begin{flushleft}
    \footnotesize
    $^\dagger$ Raw Drake estimate: $8.8\times10^6$; capped at $10^4$
    for computational tractability (see Section \ref{sec:limitations}).\\
    $^\ddagger$ Numerical underflow in \textsc{scipy} \texttt{float64};
    clustering is more extreme than double precision can represent.
  \end{flushleft}
\end{table*}

Under pessimistic parameters the destruction rate is $16.1 \pm 1.8$
per cent (coefficient of variation CV~=~11.4 per cent across seeds),
confirming that the result is robust to initialisation.
Under optimistic parameters the rate rises to $85.8 \pm 0.2$ per cent
(CV~=~0.26 per cent), a factor of~5 increase that reflects the
dramatically higher encounter rate at $N_c = 10{,}000$.

\subsection{Spatial clustering and the obs/Poisson ratio}
\label{sec:clustering}

For each seed and scenario we compute the mean nearest-neighbour
distance $\bar{d}_\text{nn}$ among surviving active civilisations and
compare it to the Poisson expectation for a homogeneous random field
of the same density:
\begin{equation}
  d_\text{Poisson} = 0.5538 \left(\frac{V_\text{eff}}{N_\text{active}}\right)^{1/3}
  \label{eq:poisson_nn}
\end{equation}
where $V_\text{eff}$ is the effective galactic volume sampled.
The ratio $r_\text{obs/P} = \bar{d}_\text{nn}/d_\text{Poisson}$
quantifies departure from randomness: values $>1$ indicate that
survivors are more dispersed than a Poisson field (super-Poissonian),
while values $<1$ indicate clustering (sub-Poissonian).

\subsubsection{Pessimistic regime: super-Poissonian survivors}

In the pessimistic scenario, $r_\text{obs/P} = 1.42 \pm 0.05$
across 10 seeds.
A KS test against the theoretical Poisson nearest-neighbour
distribution yields $p = 1.91\times10^{-42}$, confirming highly
significant departure.
The interpretation is that Dark Forest dynamics act as a
\emph{spatial filter}: when density is low, the survivors are
preferentially those that were already isolated before the first
detection event, creating voids of size $\sim 1800$\,ly around each
survivor.
This is qualitatively consistent with the void-based survival model
of \citet{forgan2011}.

\subsubsection{Optimistic regime: sub-Poissonian survivors,
  a previously unreported inversion}
\label{sec:inversion}

In the optimistic scenario, $r_\text{obs/P} = 0.578 \pm 0.020$,
a statistically unambiguous inversion of the clustering regime.
The KS $p$-value is below \textsc{float64} precision
($p < 2.2\times10^{-308}$).

This inversion arises because at high civilisation density
($N_c = 10{,}000$ in a disc of radius $\sim 50{,}000$\,ly),
the GHZ density gradient dominates over the Dark Forest
dispersal effect.
Survivors are overwhelmingly concentrated in the GHZ ring
($r_\text{peak} = 25{,}000$\,ly), producing a sub-Poissonian
distribution relative to the uniform-density Poisson expectation.
The spatial structure of survivors transitions from
\emph{void-dominated} (sparse: survivors in the gaps between
destruction events) to \emph{GHZ-dominated} (dense: survivors
clustered in the habitability sweet spot).

This result was not reported by \citet{forgan2009},
\citet{forgan2011}, or \citet{foucher2023}.
\citet{foucher2023} reported sub-Poissonian distributions using
the dispersion index $R = \sigma^2/\mu$ and the pair correlation
function $g(r)$, attributing them to predator--prey dynamics
rather than the GHZ structural effect identified here.

\subsection{Sobol sensitivity analysis}
\label{sec:sobol_results}

Table \ref{tab:sobol} reports first-order ($S_1$) and
total-effect ($S_T$) Sobol indices for the four parameters.

\begin{table}
  \centering
  \caption{Sobol sensitivity indices for the destruction rate.
           $N = 512$ Saltelli samples (5{,}120 evaluations).}
  \label{tab:sobol}
  \begin{tabular}{lcccc}
    \toprule
    Parameter & Range & $S_1$ & $S_T$ \\
    \midrule
    $\theta_\text{threat}$ & $[0.3,0.9]$     & $0.122\pm0.115$ & $0.934\pm0.115$ \\
    $\alpha$               & $[0.001,0.1]$   & $0.103\pm0.108$ & $1.026\pm0.153^\star$ \\
    $f_\text{aggressive}$  & $[0.1,0.5]$     & $0.091\pm0.105$ & $0.932\pm0.137$ \\
    $r_\text{max}$         & $[50,500]$\,ly  & $0.112\pm0.119$ & $1.072\pm0.138^\star$ \\
    \bottomrule
  \end{tabular}
  \begin{flushleft}
    \footnotesize
    $^\star$ Marginally exceeds unity due to Monte Carlo noise in the
    Saltelli estimator for stochastic outputs with CV\,=\,72\,per cent;
    both values lie within $1\sigma$ of 1.0 \citep{saltelli2010}.
  \end{flushleft}
\end{table}

All first-order indices are similar ($S_1 \approx 0.09$--$0.12$),
indicating no single dominant parameter.
Total-effect indices are uniformly large ($S_T \approx 0.93$--$1.07$),
indicating strong high-order interactions.

% ─────────────────────────────────────────────────────────────────────────────
\section{Discussion}
\label{sec:discussion}
% ─────────────────────────────────────────────────────────────────────────────

\subsection{The density-dependent regime transition}

The central finding is a density-dependent inversion in survivor
spatial structure from super-Poissonian (sparse) to sub-Poissonian
(dense).
In the sparse regime, the Dark Forest \emph{dispersal effect}
dominates: survivors tend to be initially isolated, biasing their
distribution toward voids and raising $r_\text{obs/P} > 1$.
In the dense regime, the \emph{GHZ concentration effect} dominates:
when encounters are ubiquitous, survival probability is roughly
uniform and the survivor distribution mirrors the GHZ birth
distribution, which is concentrated relative to the uniform Poisson
baseline, giving $r_\text{obs/P} < 1$.

The transition occurs near the density where the mean
inter-civilisation distance equals $r_\text{max}$.
In the pessimistic scenario the mean distance is $\sim 2{,}200$\,ly
$> r_\text{max}$; in the optimistic scenario it is $\sim 800$\,ly
$< r_\text{max}$.

\subsection{Comparison with previous work}

\citet{forgan2009} found $\sim 361$ civilisations under ``Rare Life''
parameters with $P_\text{destroy} = 0.5$, comparable in scale to
our pessimistic scenario.
His model did not compute spatial statistics of survivors.
The present model extends his framework with a physically motivated
detection function, sociological heterogeneity, and explicit spatial
analysis.

\citet{foucher2023} reported sub-Poissonian survivors in high-density
ecological models, consistent with our optimistic scenario.
The GHZ-concentration mechanism identified here is complementary to
the predator--prey minimum-distance mechanism of \citet{foucher2023}.

\subsection{The attack-threshold free parameter}

No published study constrains $\theta_\text{threat}$ empirically
\citep{devladar2013}.
The Sobol analysis shows $S_1(\theta_\text{threat}) = 0.122$
(largest first-order contributor), but the qualitative inversion
result is robust across the full $[0.3,\,0.9]$ range examined.
Bayesian inference against the SETI non-detection record could
provide informative priors and is a natural extension.

% ─────────────────────────────────────────────────────────────────────────────
\section{Limitations and future work}
\label{sec:limitations}
% ─────────────────────────────────────────────────────────────────────────────

\textit{Computational cap.}
The raw optimistic Drake estimate is $N_c = 8.8\times10^6$,
capped at $10^4$ (0.11 per cent) for tractability.
The cap affects absolute rates but not the qualitative inversion,
which depends on the ratio of mean separation to $r_\text{max}$.
Vectorising the pair-detection loop with \textsc{numpy} broadcasting
would enable $N_c \sim 10^6$ and is left for future work.

\textit{Stellar proper motion.}
Fixed positions ignore stellar drift \citep[$\sim 30$\,km\,s$^{-1}$;][]{carrollnellenback2019},
which would diffuse the GHZ concentration and reduce the magnitude
of the sub-Poissonian signal.

\textit{Free parameters.}
Bayesian inference against SETI observations would replace ad-hoc
values of $\theta_\text{threat}$, sociological fractions, and
detection noise rates with data-driven posteriors.

% ─────────────────────────────────────────────────────────────────────────────
\section{Conclusions}
\label{sec:conclusions}
% ─────────────────────────────────────────────────────────────────────────────

\begin{enumerate}
  \item Under pessimistic Drake parameters ($N_c = 440$, 10 seeds),
        the Dark Forest mechanism destroys $16.1 \pm 1.8$\,per cent
        of civilisations, producing a super-Poissonian survivor
        distribution ($r_\text{obs/P} = 1.42\pm0.05$,
        KS $p = 1.91\times10^{-42}$, CV\,=\,11.4\,per cent).

  \item Under optimistic parameters ($N_c = 10{,}000$, 10 seeds),
        the rate rises to $85.8 \pm 0.2$\,per cent and the
        survivor distribution \emph{inverts} to sub-Poissonian
        ($r_\text{obs/P} = 0.578\pm0.020$,
        KS $p < 2.2\times10^{-308}$, CV\,=\,0.26\,per cent).

  \item This density-dependent inversion of the obs/Poisson ratio
        has not been previously reported and constitutes the
        primary novel contribution.
        It arises from competition between a Dark Forest dispersal
        effect (dominant at low density) and a GHZ concentration
        effect (dominant at high density).

  \item Global Sobol analysis ($N=512$, 5{,}120 evaluations)
        confirms that no single parameter dominates
        ($S_1 \in [0.09,\,0.12]$) and that results are robust
        to parameter uncertainty.

  \item The attack threshold $\theta_\text{threat} = 0.7$ is
        the largest first-order Sobol contributor ($S_1 = 0.122$)
        but remains observationally unconstrained; Bayesian
        inference against the SETI non-detection record is
        recommended as an immediate next step.
\end{enumerate}

Code and data will be deposited at
\url{https://github.com/juan-galaz/dark-forest-mc}
upon acceptance.

\section*{Acknowledgements}

The author thanks the open-source communities behind
\textsc{numpy}, \textsc{scipy}, \textsc{pandas}, \textsc{SALib}
and \textsc{pytest}.
This research received no external funding.

\section*{Data Availability}

All result files (\texttt{multiseed\_results.json},
\texttt{multiseed\_optimista\_results.json},
\texttt{sobol\_bosque\_oscuro.json},
\texttt{resumen\_analisis.json}) and the regression test suite
(5/5 passing, seeds 42, 7, 13, 99, 256) will be deposited
in a public repository upon acceptance.

\appendix
\section{Model parameters}
\label{app:params}

\begin{table}
  \centering
  \caption{Drake filter parameters for both scenarios.}
  \label{tab:scenarios}
  \begin{tabular}{lccc}
    \toprule
    Scenario & $f_l$ & $f_i$ & $f_c$ \\
    \midrule
    Pessimistic & $10^{-3}$ & $10^{-3}$ & $10^{-2}$ \\
    Optimistic  & $0.10$    & $0.01$    & $0.20$    \\
    \bottomrule
  \end{tabular}
\end{table}

\begin{table}
  \centering
  \caption{Fixed model parameters.}
  \label{tab:params}
  \begin{tabular}{lll}
    \toprule
    Parameter & Value & Reference \\
    \midrule
    $N_\star$ & $2\times10^{11}$ & -- \\
    $r_\text{max}$ & 1000\,ly & \citet{troitskij1989,loeb2007} \\
    $r_\text{peak}$ (GHZ) & 25{,}000\,ly & \citet{lineweaver2004} \\
    $\sigma_r$ (GHZ) & 8{,}000\,ly & \citet{lineweaver2004} \\
    $L_\text{ext}$ & $1.43\times10^8$\,yr & \citet{sandberg2018} \\
    $\theta_\text{threat}$ & 0.70 & free \\
    $\alpha$ & 0.03\,yr$^{-1}$ & free \\
    $p_\text{fn}$ & 0.05 & free \\
    $p_\text{fp}$ & 0.01 & free \\
    Sobol $N$ & 512 & \citet{saltelli2010} \\
    \bottomrule
  \end{tabular}
\end{table}

\begin{thebibliography}{99}

\bibitem[\protect\citeauthoryear{Brin}{1983}]{brin1983}
Brin G.~D., 1983, QJRAS, 24, 283

\bibitem[\protect\citeauthoryear{Carroll-Nellenback et al.}{2019}]{carrollnellenback2019}
Carroll-Nellenback J., Frank A., Wright J., Scharf C., 2019, AJ, 158, 117.
\newblock \doi{10.3847/1538-3881/ab31a3}

\bibitem[\protect\citeauthoryear{Cixin}{2008}]{cixin2008}
Cixin L., 2008, The Dark Forest. Chongqing Publishing House

\bibitem[\protect\citeauthoryear{de Vladar}{2013}]{devladar2013}
de Vladar H.~P., 2013, IJA, 12, 53.
\newblock \doi{10.1017/S1473550412000407}

\bibitem[\protect\citeauthoryear{Fermi}{1950}]{fermi1950}
Fermi E., 1950, unpublished remarks, Los Alamos

\bibitem[\protect\citeauthoryear{Forgan}{2009}]{forgan2009}
Forgan D.~H., 2009, IJA, 8, 121.
\newblock \doi{10.1017/S1473550408004321}

\bibitem[\protect\citeauthoryear{Forgan}{2011}]{forgan2011}
Forgan D.~H., 2011, IJA, 10, 341.
\newblock \doi{10.1017/S1473550411000127}

\bibitem[\protect\citeauthoryear{Foucher}{2023}]{foucher2023}
Foucher F., 2023, AcAau, 204, 785.
\newblock \doi{10.1016/j.actaastro.2022.12.007}

\bibitem[\protect\citeauthoryear{Hart}{1975}]{hart1975}
Hart M.~H., 1975, QJRAS, 16, 128

\bibitem[\protect\citeauthoryear{Lineweaver, Fenner \& Gibson}{2004}]{lineweaver2004}
Lineweaver C.~H., Fenner Y., Gibson B.~K., 2004, Science, 303, 59.
\newblock \doi{10.1126/science.1092322}

\bibitem[\protect\citeauthoryear{Loeb \& Zaldarriaga}{2007}]{loeb2007}
Loeb A., Zaldarriaga M., 2007, JCAP, 2007, 020.
\newblock \doi{10.1088/1475-7516/2007/01/020}

\bibitem[\protect\citeauthoryear{Maneewongvatana \& Mount}{1999}]{maneewongvatana1999}
Maneewongvatana S., Mount D.~M., 1999, preprint (arXiv:cs/9901013)

\bibitem[\protect\citeauthoryear{Saltelli et al.}{2010}]{saltelli2010}
Saltelli A. et al., 2010, Comput.~Phys.~Commun., 181, 259.
\newblock \doi{10.1016/j.cpc.2009.09.018}

\bibitem[\protect\citeauthoryear{Sandberg, Drexler \& Ord}{2018}]{sandberg2018}
Sandberg A., Drexler E., Ord T., 2018, AcAau, 150, 24.
\newblock \doi{10.1016/j.actaastro.2018.04.030}

\bibitem[\protect\citeauthoryear{Troitskij}{1989}]{troitskij1989}
Troitskij V.~S., 1989, AcAau, 19, 875.
\newblock \doi{10.1016/0094-5765(89)90079-9}

\bibitem[\protect\citeauthoryear{Webb}{2002}]{webb2002}
Webb S., 2002, If the Universe Is Teeming with Aliens. Copernicus Books, New York

\end{thebibliography}

\end{document}
